\section{Uncertainty Quantification Applications}
In this section, we will discuss several important applications that can be tackled using our proposed post-hoc uncertainty learning framework. 
%Given the pretrained model $\vh$ and its training data $\mathcal{D}_{B}=\left\{\boldsymbol{x}^B_{i}, y^B_{i}\right\}_{i=1}^{N_B}$, the goal is to efficiently train the meta-model $\vg$ using the exiting training data. 
Our meta-model has two major advantages: (1) It is efficient, i.e., it is easy to optimize and empirically require less number of training data due to smaller model complexity, (2) It is versatile, i.e., it not only can quantify different uncertainties to adapt to different applications, but also different settings, i.e., like transfer learning, the distribution of training data for meta-model can even be different from the base-model.
\paragraph{Out of domain data detection}
The base-model $\vh$ is trained using data sampled from distribution $P^B_{Z}$. During testing, if there exists some unobserved out of domain data from another distribution $P^{ood}_{Z}$, the meta-model should correctly identify the in-distribution or out-of-distribution input sample. We use the same training set as base model to train the meta-model, i.e., $\mathcal{D}_{B} = \mathcal{D}_{M}$, which means our proposed method does not require additional OOD data for training. For OOD task, we use three epistemic uncertainties to measure the uncertainty score $\vu\left(\tilde{y}\right)$.
\paragraph{Miss-classification Detection}
Instead of detecting whether a testing sample is out of domain, the meta-model aims to identify the failure or success of prediction during testing. We train the meta-model with the same setting as OOD detection task, and use the two total uncertainties.
\paragraph{Trustworthy Transfer Learning}
The aforementioned applications assume the data used for training the base-model and the meta-model are exactly same, i.e., $\mathcal{D}_{B} = \mathcal{D}_{M}$, or $P^B_{Z} = P^B_{M}$ in general. Instead, the transfer learning assumes there exists a pretrained model $\vh_s$ trained using source task data $\mathcal{D}_{s}$ sampled from source distribution $P^s_{Z}$, and the goal is to adapt the source model to target task using target data $\mathcal{D}_{t}$ sampled from target distribution $P^t_{Z}$. Most existing transfer learning problem focus on improving the prediction performance of transferred model, but ignore its uncertainty quantification performance on the target task. Our meta-model method can be utilized as a simple and effective way to address this problem, i.e, given pretrained source model $\vh \circ \vf = \argmin_{\vh \circ \vf} \gL_E(\vh \circ \vf,\gD_s)$, the meta-model can be efficiently trained using target domain data, i.e., $\vg = \argmin_{\vg} \gL_E(\vg \circ \vf,\gD_t)$.  